% !TEX program = lualatex
\documentclass[letterpaper,twoside]{book}

% ==================================================
%      COMPILAR CON XeLaTeX o LuaLaTeX
% ==================================================
\usepackage{changepage}
\usepackage{float}
\usepackage{multicol}

% --- Idioma y codificación ---
\usepackage{polyglossia}
\setmainlanguage{spanish}

% --- Tipografía ---
\usepackage{fontspec}
\setmainfont{Latin Modern Roman}
\usepackage{microtype}

% --- Matemáticas ---
\usepackage{amsmath,amssymb,amsthm,amsfonts}
\usepackage{mathtools}

% --- Márgenes para impresión doble cara ---
\usepackage[
  top=2.5cm,
  bottom=2.5cm,
  inner=3cm,
  outer=2.5cm,
  headheight=14pt,
  headsep=10pt
]{geometry}

% --- Corrección de páginas en blanco impares (impresión) ---
\makeatletter
\let\origcleardoublepage\cleardoublepage
\def\cleardoublepage{%
  \clearpage%
  \ifodd\c@page\else
    \hbox{}%
    \thispagestyle{empty}%
    \newpage%
  \fi}
\makeatother

% --- Colores personalizados ---
\usepackage{xcolor}
\definecolor{dlblue}{RGB}{31,73,125}
\definecolor{dlblueLight}{RGB}{230,239,250}
\definecolor{dlgreen}{RGB}{22,101,52}
\definecolor{dlgreenLight}{RGB}{232,246,237}
\definecolor{dlred}{RGB}{165,29,45}
\definecolor{dlorange}{RGB}{146,64,14}
\definecolor{dlorangeLight}{RGB}{255,244,229}
\definecolor{dlpurple}{RGB}{128,0,128}
\definecolor{dlcodebg}{RGB}{247,248,250}
\definecolor{commentgreen}{RGB}{0,128,0}
\definecolor{keywordblue}{RGB}{0,0,200}
\definecolor{stringred}{RGB}{163,21,21}

% --- Cajas personalizadas ---
\usepackage[most]{tcolorbox}

% --- Contadores por capítulo ---
\newcounter{definicion}[chapter]
\renewcommand{\thedefinicion}{\thechapter.\arabic{definicion}}

\newcounter{algoritmo}[chapter]
\renewcommand{\thealgoritmo}{\thechapter.\arabic{algoritmo}}

\newcounter{advertencia}[chapter]
\renewcommand{\theadvertencia}{\thechapter.\arabic{advertencia}}

\newcounter{ejercicio}[chapter]
\renewcommand{\theejercicio}{\thechapter.\arabic{ejercicio}}

% --- Configuración Global de Cajas ---
\tcbset{
  box_style/.style={
    enhanced,
    colback=white,
    boxrule=0.6pt,
    left=6pt, right=6pt, top=6pt, bottom=6pt,
    sharp corners,
    before skip=12pt,
    after skip=12pt,
    fonttitle=\bfseries,
    attach title to upper,
    after title={\par\nobreak\vskip4pt\hrule\vskip6pt\noindent},
  }
}

% --- Definición (azul) ---
\newtcolorbox{definicionBox}[1][]{
  box_style,
  colframe=dlblue,
  coltitle=dlblue,
  borderline west={2pt}{0pt}{dlblue},
  title={Definición \thedefinicion: #1},
}
\newenvironment{definicion}[1][]{%
  \refstepcounter{definicion}\begin{definicionBox}[#1]\normalfont}{\end{definicionBox}}

% --- Algoritmo (púrpura) ---
\newtcolorbox{algoritmoBox}[1][]{
  box_style,
  colframe=dlpurple,
  coltitle=dlpurple,
  borderline west={2pt}{0pt}{dlpurple},
  title={Algoritmo \thealgoritmo: #1},
}
\newenvironment{algoritmo}[1][]{%
  \refstepcounter{algoritmo}\begin{algoritmoBox}[#1]\normalfont}{\end{algoritmoBox}}

% --- Advertencia (rojo) ---
\newtcolorbox{advertenciaBox}[1][]{
  box_style,
  colframe=dlred,
  coltitle=dlred,
  borderline west={2pt}{0pt}{dlred},
  title={Advertencia \theadvertencia: #1},
}
\newenvironment{advertencia}[1][]{%
  \refstepcounter{advertencia}\begin{advertenciaBox}[#1]\normalfont}{\end{advertenciaBox}}

% --- Ejercicio (sin borde, solo texto) ---
\newenvironment{ejercicio}
  {\par\medskip\noindent\refstepcounter{ejercicio}\textbf{Ejercicio \theejercicio.}}
  {\par\medskip}

% --- Código fuente ---
\usepackage{listings}

\lstdefinelanguage{PythonIC}{
  morekeywords={import,from,as,def,return,if,else,for,while,print,class,None,True,False,in,range,len,float,argmax,reshape},
  morekeywords=[2]{keras,np,tf,plt,red,imgs,labs,X,Y,h,fig,axs,model},
  sensitive=true,
  morecomment=[l]{\#},
  morestring=[b]',
  morestring=[b]",
}

\lstset{
  language=PythonIC,
  basicstyle=\ttfamily\footnotesize,
  keywordstyle=\color{keywordblue}\bfseries,
  keywordstyle=[2]\color{dlgreen},
  commentstyle=\color{commentgreen}\itshape,
  stringstyle=\color{stringred},
  backgroundcolor=\color{dlcodebg},
  frame=single,
  rulecolor=\color{gray!40},
  tabsize=2,
  showstringspaces=false,
  breaklines=true,
  breakatwhitespace=false,
  breakautoindent=false,
  postbreak=\mbox{\hbox{\color{gray}\tiny$\hookrightarrow$}\space},
  numbers=left,
  numberstyle=\tiny\color{gray},
  numbersep=8pt,
  xleftmargin=15pt,
  framexleftmargin=10pt,
  captionpos=b,
}

% --- Diagramas y gráficas ---
\usepackage{tikz}
\usetikzlibrary{positioning,arrows.meta,shapes.geometric,calc,fit}

% --- Tablas profesionales ---
\usepackage{booktabs}
\usepackage{tabularx}

% --- Listas personalizadas ---
\usepackage{enumitem}

% --- Cabeceras y pies de página ---
\usepackage{fancyhdr}
\usepackage{lastpage}

% --- Evitar líneas desbordadas en párrafos normales ---
\emergencystretch=2em

% ==================================================
%   CABECERAS Y PIES DE PÁGINA
% ==================================================
\pagestyle{fancy}
\fancyhf{}
\fancyhead[LE]{\small\textit{\leftmark}}
\fancyhead[RO]{\small\textit{\rightmark}}
\fancyfoot[LE,RO]{\small\thepage\ de \pageref{LastPage}}
\fancyfoot[CE,CO]{\small\textit{Deep Learning --- DLY0100}}
\renewcommand{\headrulewidth}{0.4pt}
\renewcommand{\footrulewidth}{0.4pt}

% Sin headers/footers en preface e índice
\fancypagestyle{plain}{%
  \fancyhf{}%
  \fancyfoot[C]{\small\thepage}%
  \renewcommand{\headrulewidth}{0pt}%
  \renewcommand{\footrulewidth}{0pt}%
}

% Título de las partes y capítulos
\renewcommand{\partname}{Parte}
\renewcommand{\thepart}{\Roman{part}}

% ==================================================
%   HIPERVÍNCULOS Y METADATOS DEL DOCUMENTO
% ==================================================
\usepackage{hyperref}
\hypersetup{
  pdftitle={Un Apunte de Deep Learning},
  pdfauthor={Antonio Jara Sánchez},
  colorlinks=true,
  linkcolor=black,
  urlcolor=black,
  citecolor=black,
}

% ==================================================
%   COMIENZO DEL DOCUMENTO
% ==================================================
\begin{document}

\thispagestyle{empty}

% --- Portada ---
\begin{titlepage}
\begin{center}
\vspace*{\fill}

% Línea horizontal superior
\rule{\textwidth}{1pt}
\vspace{1cm}

{\Large Un Apunte de}

\vspace{0.5cm}
{\Huge \textbf{Deep Learning}}

\vspace{1cm}
\rule{\textwidth}{0.8pt}

\vspace{2.5cm}

\Large \textit{Antonio Jara Sánchez}

\vspace{3cm}

\normalsize Versión: \today

\vspace*{\fill}
\end{center}
\end{titlepage}
\cleardoublepage

% --- Prefacio ---
\pagenumbering{roman}
\pagestyle{plain}
\chapter*{Prefacio}
\addcontentsline{toc}{chapter}{Prefacio}

Este apunte, como el intento de presentar los temas del curso de forma progresiva que es, fue realizado para fines de estudio personal de la asignatura \textbf{DLY0100}. Se basa casi en su totalidad en el material del curso (diapositivas), así como notas de clase.

\medskip

Esta es la versión del \today. Cualquier sugerencia, corrección, observación, etc. puede ser indicado al correo antoniojarasanchez@proton.me

\cleardoublepage

% --- Índice ---
\tableofcontents

\cleardoublepage
\pagenumbering{arabic}
\setcounter{page}{1}
\pagestyle{fancy}

% --- Inicio del contenido ---
% ==================================================
%  CAPÍTULO 1: INTRODUCCIÓN AL DEEP LEARNING
% ==================================================

\section{Introducción al Deep Learning}

\subsection{Contexto tecnológico: la Sociedad 5.0}

En 2015, Japón acuñó el concepto de \textbf{Sociedad 5.0}, cuya idea central es que \emph{las máquinas serán las principales creadoras del conocimiento humano}, apalancadas en la inteligencia artificial. Mientras la Sociedad 4.0 generó la información digital, la 5.0 la aprovecha con IA para integrar la tecnología en la vida cotidiana.

\subsection{Inteligencia Artificial: definición y evolución}

Se adopta la definición de Nilsson (1998):

\begin{definicion}
\textbf{Inteligencia Artificial (IA)}: capacidad de imitar comportamientos inteligentes.
\end{definicion}

La definición de IA ha evolucionado a lo largo del tiempo.

\begin{table}[H]
\centering
\caption{Evolución de la definición de Inteligencia Artificial.}
\label{tab:evolucion-ia}
\small
\begin{tabularx}{\textwidth}{@{}lXX@{}}
\toprule
\textbf{Definición histórica} & \textbf{Idea central} & \textbf{Problema / observación} \\
\midrule
Sistemas que piensan como humanos & Imitar el pensamiento humano & No existe una única forma de pensamiento humano \\
Sistemas que piensan racionalmente & Pensar de manera correcta & El silogismo carece de contexto \\
Sistemas que actúan como humanos & Test de Turing & No tuvo éxito en su época por capacidad de procesamiento \\
Sistemas que actúan racionalmente & Actuar según lo programado & Definición vigente desde 1998 \\
\bottomrule
\end{tabularx}
\end{table}

\begin{definicion}
\textbf{IA (definición vigente)}: todo sistema capaz de \emph{actuar racionalmente}.

\textbf{Racional}: correcto para lo que fue programado; no implica conciencia.
\end{definicion}

\subsubsection{Clasificaciones de la Inteligencia Artificial}

\begin{table}[H]
\centering
\caption{Clasificación de la IA según capacidad y complejidad.}
\label{tab:ia-capacidad}
\small
\begin{tabularx}{\textwidth}{@{}lXX@{}}
\toprule
\textbf{Tipo} & \textbf{Descripción} & \textbf{Ejemplos} \\
\midrule
IA reactiva & No aprende de experiencias pasadas; responde a estímulos actuales & Alexa, Siri, aspiradora robot, Deep Blue \\
IA con memoria limitada & Aprende de datos históricos, pero con memoria limitada y específica & Vehículos autónomos, Waze, Google Maps \\
\bottomrule
\end{tabularx}
\end{table}

\begin{table}[H]
\centering
\caption{Clasificación de la IA según propósito y aplicaciones.}
\label{tab:ia-proposito}
\small
\begin{tabularx}{\textwidth}{@{}lXX@{}}
\toprule
\textbf{Tipo} & \textbf{Descripción} & \textbf{Ejemplos} \\
\midrule
IA débil / narrow & Específica para tareas claras; no aprende tareas nuevas & Detector de spam de Gmail, asistentes virtuales, sistemas de recomendación, reconocimiento de voz \\
IA fuerte & Capaz de aprender cosas nuevas & Autos autónomos, sistemas de Machine Learning \\
Super IA & Consciente de su existencia; capaz de realizar cualquier tarea humana & No existe todavía \\
\bottomrule
\end{tabularx}
\end{table}

\subsection{Machine Learning: el núcleo de la IA}

\begin{definicion}
\textbf{Machine Learning} (aprendizaje automático): campo de la IA que se encarga de reconocer
patrones desde los datos y de generalizar conocimiento a partir de la experiencia y de muchos
datos.

Definición de \textbf{Arthur Samuel}: ``el diseño y desarrollo de algoritmos que permiten a los
computadores \emph{aprender sin ser explícitamente programados}''.
\end{definicion}

No existe un comando ``\texttt{.learn}'', el aprendizaje se logra mediante un \textbf{algoritmo}.

\begin{definicion}
\textbf{Algoritmo}: secuencia lógica, paso a paso, con inicio, decisiones y fin u objetivo.
\end{definicion}

\begin{table}[H]
\centering
\caption{Tipos de aprendizaje automático.}
\label{tab:tipos-ml}
\small
\begin{tabularx}{\textwidth}{@{}lXX@{}}
\toprule
\textbf{Tipo} & \textbf{Descripción} & \textbf{Ejemplos} \\
\midrule
Supervisado & Aprende con ejemplos etiquetados & Detección de spam, reconocimiento de voz, analíticas de anuncios \\
No supervisado & Encuentra patrones sin etiquetas previas & Sistemas de recomendación, registros de conductas, detección de anomalías \\
Semi-supervisado & Mezcla de supervisado y no supervisado & Análisis de voz, detección de fraude, investigaciones médicas \\
Por refuerzo & Aprende mediante prueba, error y refuerzo & Robótica, videojuegos, gestión de recursos, IoT \\
\bottomrule
\end{tabularx}
\end{table}

\subsection{Deep Learning: aprendizaje por capas}

\begin{definicion}
\textbf{Deep Learning} (aprendizaje profundo): sistema capaz de aprender mediante algoritmos y
\emph{sin intervención manual} de un usuario.
\end{definicion}

\begin{definicion}
\textbf{Representación}: forma en que los datos son transformados y procesados a lo largo de
las capas para extraer características relevantes de manera automática.

\textbf{Profundidad del modelo}: cantidad de capas sucesivas de representación.
\end{definicion}

\begin{advertencia}
``Profundo'' \emph{no} significa comprensión más profunda; significa \emph{muchas capas
sucesivas de representación}.
\end{advertencia}

\textbf{¿Por qué despegó después de 2012?}

Las ideas clave de la visión artificial (redes neuronales convolucionales y retropropagación)
ya se entendían en 1990 y el algoritmo \textbf{LSTM} (memoria a corto y largo plazo) se
desarrolló en 1997 y apenas ha cambiado. Pero hasta 2010, un computador de un terabyte era muy difícil
de conseguir.

\textbf{Tres nuevas figuras han impulsado el avance del aprendizaje automático:}
\begin{enumerate}[topsep=2pt,itemsep=1pt]
  \item Hardware más potente (GPU, TPU).
  \item Mayor disponibilidad de datos (big data).
  \item Avances en los algoritmos (mejores funciones de activación, optimizadores).
\end{enumerate}

\subsubsection{Machine Learning vs Deep Learning}

\begin{table}[H]
\centering
\caption{Comparación entre Machine Learning y Deep Learning.}
\label{tab:ml-vs-dl}
\small
\begin{tabularx}{\textwidth}{@{}lXX@{}}
\toprule
\textbf{Criterio} & \textbf{Machine Learning} & \textbf{Deep Learning} \\
\midrule
Extracción de variables & Manual, mediante ingeniería de características & Automática, mediante capas de representación \\
Datos & Principalmente estructurados & Imágenes, texto, voz y otros no estructurados \\
Intervención manual & Requiere decirle si aprendió bien o mal & Aprende sin intervención manual \\
Escalabilidad & Computacionalmente eficiente, pero no se adapta bien a grandes volúmenes & Suele ser más eficaz en problemas complejos y grandes volúmenes \\
Costo computacional & Menor & Mayor \\
Interpretabilidad & Mayor & Menor \\
\bottomrule
\end{tabularx}
\end{table}

\begin{advertencia}
El Deep Learning es más caro y menos interpretable, pero más eficaz en problemas complejos.
\end{advertencia}
% ==================================================
%  CAPÍTULO 2: EL PERCEPTRÓN Y REDES FULLY CONNECTED
% ==================================================
\cleardoublepage
\section{El Perceptrón y Redes Fully Connected Feed Forward}

\subsection{De la IA a las redes neuronales}

Una red se puede pensar como una malla de nodos entrelazados con diferentes uniones. En Deep Learning programamos un sistema que en cada capa decide hacia qué nodo avanzar según una probabilidad dada por un algoritmo.

En \textbf{\emph{Representation Learning}} (aprendizaje de representaciones) se construyen representaciones de las entradas o inputs en etapas sucesivas. En cada etapa, la representación o abstracción de la anterior contiene información más compacta y relevante, pero en menor cantidad. Eventualmente, con esta abstracción sucesiva se crea un tensor.Porque transforman la entrada inicial hasta generar un tensor o vector final que contiene todo lo necesario para tomar una decisión, estos algoritmos son llamados \textbf{\emph{End-to-End}}.

La diferencia fundamental entre Machine Learning y Deep Learning radica en la \emph{extracción de características}:

\begin{table}[H]
\centering
\caption{Extracción de características: ML vs DL.}
\label{tab:representacion}
\small
\begin{tabular}{@{}ll@{}}
\toprule
\textbf{Enfoque} & \textbf{Proceso} \\
\midrule
Machine Learning & Extracción manual de características $\to$ clasificador $\to$ salida \\
Deep Learning & Entrada $\to$ red neuronal $\to$ salida, sin intervención manual \\
\bottomrule
\end{tabular}
\end{table}

\paragraph{Rendimiento, datos y costo computacional:}
A diferencia de los modelos tradicionales de Machine Learning (cuyo rendimiento se satura rápidamente al aumentar los datos), los modelos de Deep Learning escalan a medida que se incrementan la cantidad de datos, el número de capas y la capacidad computacional.
La idea es que se \emph{podría} resolver prácticamente cualquier problema, dados una cantidad de datos y capacidad de cómputo suficientes, siendo estos mismos el factor limitante.

\subsection{El perceptrón}

El \textbf{perceptrón} se inspira en el funcionamiento de la red neuronal biológica del cerebro humano. En este, cada neurona recibe impulsos eléctricos de neuronas vecinas. Cuando la cantidad total de las señales entrantes supera cierto \textbf{umbral}, la neurona ``dispara'' su propia señal, la cual se conecta con otras neuronas.

La unión o punto de contacto entre dos neuronas se llama \textbf{sinapsis}. Cuando aprendemos nuevas asociaciones entre dos conceptos, la \textbf{fuerza sináptica} entre las neuronas que representan dichos conceptos se fortalece.

\textbf{Regla de Hebb (1949)}: ``Las células que se \textbf{activan juntas}, se conectan entre sí''.

\begin{definicion}
\textbf{Perceptrón (Rosenblatt, 1957)}: modelo simplificado de una neurona biológica llevado
al computador.
\end{definicion}
\end{document}
